\section{Future Scopes}

This project already shows that static recommendation is not enough when user preference keeps changing. 
At the same time, the results also showed a few clear limits in the current design. 
The future work of this system is therefore not about adding random new ideas. 
It is about fixing the missing pieces that matter most.



%%%%%%%%%%%%%%%%%%%%%%%%%%%%%%%%%%%%%%%%%%%%%%%%%%%%%%%%%%%%%%%%
%%%%%%%%%%%%%%%%%%%%%%%%%%%%%%%%%%%%%%%%%%%%%%%%%%%%%%%%%%%%%%%%

\subsection{Gating Mechanism}

The biggest missing module is the gating mechanism.

Right now, the system has three useful modes:

\begin{itemize}
    \item \textbf{S1} for fast static retrieval
    \item \textbf{S2P} for updated preference learning
    \item \textbf{S2F} for cognitive-state-aware recommendation
\end{itemize}

All three are useful, but they can not be used all the time. 
S1 based embeddings quickly lose relevance as preference updates, and S2P is a very expensive module. 
The real challenge is to decide \textbf{when} to use each one. 
That is what the gating mechanism should do.
A good gate should balance all our previous goals, along with the addition of computational cost.

This is important because S2F is the strongest module, but it should not always be the default choice. 
The system becomes much more practical only when it learns to switch between the simpler and heavier modules in a controlled way.


%%%%%%%%%%%%%%%%%%%%%%%%%%%%%%%%%%%%%%%%%%%%%%%%%%%%%%%%%%%%%%%%
%%%%%%%%%%%%%%%%%%%%%%%%%%%%%%%%%%%%%%%%%%%%%%%%%%%%%%%%%%%%%%%%

\subsection{Simulation and Evaluation Limitations}

The simulation helped us test the model, but it also exposed some issues in the evaluation setup.

\subsubsection{Top-K selection bias}

In the current simulation, the next step is often driven by the first item from the top-k list. 
This makes the interaction path too deterministic. 
A random or probabilistic choice from the top-k set would give a more realistic simulation.

\subsubsection{Metric design issues}

Some metrics turned out to be weak, redundant, or hard to interpret. 
A few of them gave almost the same pattern, so they did not add much new information. 
Coverage was also tricky, because the size of the true-item set changes as the user preference changes during simulation. 
So coverage here is not just a measure of recommendation breadth. 
It is also affected by how the user state grows over time.

This shows that we need better metrics that can separate:

\begin{itemize}
    \item exploration
    \item preference growth
    \item and actual recommendation quality
\end{itemize}

\subsubsection{Missing user-level evaluation}

Most of the current evaluation is averaged over users. 
That is useful, but it hides individual failures. 
A model can look good on average even if some users fail early.

A better future evaluation should include measures like:

\begin{itemize}
    \item success rate per user
    \item how many users survive the full horizon
    \item and how many users become fatigued early
\end{itemize}


%%%%%%%%%%%%%%%%%%%%%%%%%%%%%%%%%%%%%%%%%%%%%%%%%%%%%%%%%%%%%%%%
%%%%%%%%%%%%%%%%%%%%%%%%%%%%%%%%%%%%%%%%%%%%%%%%%%%%%%%%%%%%%%%%

\subsection{UserModel Limitations}

The UserModel is useful, but it is still a handcrafted simulator.

\subsubsection{Intuition-driven behavior}

The current behavior rules were designed manually. 
They are reasonable, but they are not based on a formal psychological study. 
So the model is practical, but it is still intuitive.

\subsubsection{Tuned constants}

Several values were chosen by experiment, including:

\begin{itemize}
    \item preference volatility
    \item fatigue update weights
    \item and discontentment update weights
\end{itemize}

These constants work for the current setup, but they are still hand-tuned. 
That means the behavior is tied to our design choices, not learned from real user data. 

\subsubsection{Single-drift preference dynamics}

At the moment, preference mostly drifts toward items the user clicks. 
This works well for one kind of behavior, but not for users whose interests change in sharper or more complex ways.

\subsubsection{Multi-state preference dynamics}

A more realistic user may move between multiple preference states over time. 
For example, a user can shift between different moods, topics, or interests. 
A future simulator could sample several preference anchors and switch between them over time.

\subsubsection{Multi-base preference limitation}

The current embedding setup still assumes that a user lives in one coherent region of the space. 
But in real life, a user can like very different item groups at the same time.

It may also explain why the T0 scores are not as strong as expected. 
The same user embedding is being pulled toward different item groups, 
and that makes the final representation harder to fit cleanly.


%%%%%%%%%%%%%%%%%%%%%%%%%%%%%%%%%%%%%%%%%%%%%%%%%%%%%%%%%%%%%%%%
%%%%%%%%%%%%%%%%%%%%%%%%%%%%%%%%%%%%%%%%%%%%%%%%%%%%%%%%%%%%%%%%

\subsection{Explicit Dislike Modeling}

The current system does not track a clear dislike event. 
In practice, click and no-click are not the same as like and dislike. 
This makes the random sampling pool less controlled than it should be.

A future version should:

\begin{itemize}
    \item record disliked items explicitly
    \item avoid showing them again in random sampling
    \item and keep them out of the candidate pool when possible
\end{itemize}

This would make the simulated interaction cleaner and more realistic.



%%%%%%%%%%%%%%%%%%%%%%%%%%%%%%%%%%%%%%%%%%%%%%%%%%%%%%%%%%%%%%%%
%%%%%%%%%%%%%%%%%%%%%%%%%%%%%%%%%%%%%%%%%%%%%%%%%%%%%%%%%%%%%%%%

\subsection{Recommendation Memory}

The system does not yet keep a strict memory of what has already been recommended. 
So, in theory, the same item can appear again in the recommendation list.

That did not happen too often in our current runs, but it is still a limitation. 
A better system should track recommendation history and filter duplicates more carefully.



%%%%%%%%%%%%%%%%%%%%%%%%%%%%%%%%%%%%%%%%%%%%%%%%%%%%%%%%%%%%%%%%
%%%%%%%%%%%%%%%%%%%%%%%%%%%%%%%%%%%%%%%%%%%%%%%%%%%%%%%%%%%%%%%%

\subsection{Better Observation Space for S2F}

The current RL observation space is already useful, but it is still limited. 
Right now, S2F observes:

\begin{itemize}
    \item the updated user preference embedding
    \item recent click rate
    \item similarity of recent items
    \item diversity score
    \item and session length
\end{itemize}

This is a good start, but it may not be enough for a stronger policy. 
In future work, the observation space can include richer behavioral signals and more context from the session. 
That would give the algorithm a better view of how the user is actually reacting.



%%%%%%%%%%%%%%%%%%%%%%%%%%%%%%%%%%%%%%%%%%%%%%%%%%%%%%%%%%%%%%%%
%%%%%%%%%%%%%%%%%%%%%%%%%%%%%%%%%%%%%%%%%%%%%%%%%%%%%%%%%%%%%%%%

\subsection{Final Note}

The main value of these future directions is not that the current system is incomplete. 
It is that the current results already point to the next set of problems very clearly. 
The most important next step is to make the system smarter about \textbf{when to use each module}, 
and to make the simulation and evaluation more faithful to real user behavior.

That is where the next version of this project can become much stronger.
